\documentclass[UTF8,a4paper,12pt]{ctexart}

\usepackage[margin=2.54cm]{geometry}
\usepackage{setspace}
\usepackage{graphicx}
\usepackage[table]{xcolor}
\usepackage{tcolorbox}
\tcbuselibrary{skins}
\usepackage{booktabs}
\usepackage{array}
\usepackage{longtable}
\usepackage{tabularx}
\usepackage{float}
\usepackage{caption}
\usepackage{enumitem}
\usepackage{fancyhdr}
\usepackage{titlesec}
\usepackage{hyperref}
\usepackage{amsmath}
\usepackage{amssymb}
\usepackage[section]{placeins}
\usepackage{listings}
\usepackage{url}

\definecolor{GoogleBlue}{HTML}{1A73E8}
\definecolor{GoogleGray}{HTML}{F8F9FA}
\definecolor{GoogleBorder}{HTML}{DADCE0}
\definecolor{GoogleText}{HTML}{202124}
\definecolor{GoogleHeader}{HTML}{F1F3F4}
\definecolor{GoogleSubtle}{HTML}{5F6368}

\hypersetup{
  colorlinks=true,
  linkcolor=GoogleBlue,
  urlcolor=GoogleBlue,
  citecolor=GoogleBlue,
  pdftitle={控制平面与API接入模块实现与测试报告},
  pdfauthor={叶容宇, 胡治禾, 杨佳利}
}

\renewcommand{\familydefault}{\sfdefault}
\setstretch{1.5}
\setlength{\parindent}{2em}
\setlength{\parskip}{0.5em}
\setlength{\emergencystretch}{2em}
\sloppy
\captionsetup{font=small, labelfont={color=GoogleBlue, bf}}
\setlist[enumerate]{label=\arabic*., leftmargin=2.8em, itemsep=0.2em}
\setlist[itemize]{label=\textcolor{GoogleBlue}{\textbullet}, leftmargin=2.8em, itemsep=0.2em}
\setcounter{topnumber}{3}
\setcounter{bottomnumber}{2}
\setcounter{totalnumber}{5}
\renewcommand{\topfraction}{0.92}
\renewcommand{\bottomfraction}{0.85}
\renewcommand{\textfraction}{0.08}
\renewcommand{\floatpagefraction}{0.85}

\pagestyle{fancy}
\fancyhf{}
\cfoot{\textcolor{GoogleBorder}{\rule{0.5\textwidth}{0.4pt}}\\\vspace{0.2em}\thepage}
\renewcommand{\headrulewidth}{0pt}

\renewcommand{\thesection}{\arabic{section}}
\renewcommand{\thesubsection}{\arabic{section}.\arabic{subsection}}
\renewcommand{\thesubsubsection}{\arabic{section}.\arabic{subsection}.\arabic{subsubsection}}

\titleformat{\section}{\sffamily\zihao{3}\bfseries\color{GoogleBlue}}{\thesection}{0.5em}{}
\titleformat{\subsection}{\sffamily\zihao{4}\bfseries\color{GoogleText}}{\thesubsection}{0.5em}{}
\titleformat{\subsubsection}{\sffamily\zihao{-4}\bfseries\color{GoogleText}}{\thesubsubsection}{0.5em}{}

\titlespacing*{\section}{0pt}{1.2em}{0.6em}
\titlespacing*{\subsection}{0pt}{1em}{0.4em}
\titlespacing*{\subsubsection}{0pt}{0.8em}{0.3em}

\lstdefinestyle{asciiart}{
  basicstyle=\ttfamily\small,
  breaklines=true,
  frame=single,
  rulecolor=\color{GoogleBorder},
  backgroundcolor=\color{GoogleGray},
  columns=fullflexible,
  keepspaces=true,
  showstringspaces=false,
  xleftmargin=0.5em,
  xrightmargin=0.5em
}

\newcommand{\studentName}{叶容宇、胡治禾、杨佳利}
\newcommand{\collegeName}{计算机科学与技术学院}
\newcommand{\majorName}{软件工程}
\newcommand{\reportYear}{2026}
\newcommand{\reportMonth}{5}
\newcommand{\reportDay}{18}
\newcommand{\moduleTitle}{控制平面与 API 接入模块实现与测试报告}

\newcommand{\coverline}[2]{#1 & \underline{\makebox[8cm][c]{#2}} \\[1.1em]}
\newcommand{\docfigure}[3][0.95\textwidth]{
  \begin{figure}[!htbp]
    \centering
    \includegraphics[width=#1]{#2}
    \caption{#3}
  \end{figure}
}

\begin{document}

\begin{titlepage}
  \centering
  \vspace*{2.6cm}

  {\zihao{2}\bfseries 大规模信息系统构建技术导论\par}
  \vspace{2cm}
  {\zihao{-1}\bfseries \moduleTitle\par}

  \vfill

  \renewcommand{\arraystretch}{1.4}
  \begin{tabular}{>{\zihao{4}}p{2.8cm} >{\zihao{4}}p{9cm}}
    \coverline{姓名：}{\studentName}
    \coverline{学院：}{\collegeName}
    \coverline{专业：}{\majorName}
  \end{tabular}

  \vfill

  {\zihao{4}\reportYear~年~\reportMonth~月~\reportDay~日\par}
\end{titlepage}

\pagenumbering{arabic}
\setcounter{page}{1}

\tableofcontents
\clearpage

\zihao{-4}

\section{系统模块简介}

本人在分布式 MiniSQL 系统中负责梳理控制平面与 API 接入模块，主要涉及 \texttt{internal/controller}、\texttt{internal/apiserver}、\texttt{internal/discovery}、\texttt{internal/coordinator}、\texttt{internal/app}、\texttt{cmd/server} 与 \texttt{cmd/cli} 等目录。该模块将底层 Raft 副本组组织成可路由、可迁移、可观测的分片数据库系统，是 ShardDB 相比单副本组 DDB 的主要扩展部分。

控制平面负责维护“分片到副本组”的全局元数据，API 接入层负责把用户 SQL、控制命令和 Dashboard 请求转换为对控制平面或数据平面的调用。二者共同形成系统的接入面：用户不再直接关心某条数据位于哪一个 Raft 组，而是通过统一 API Server 或 CLI 发起请求，由系统自动完成服务发现、分片路由、Leader 选择和异常重试。

本模块的具体功能如下：
\begin{enumerate}
  \item 基于 etcd 实现服务发现，维护在线节点、节点角色、副本组、Leader 状态和 removed 标记。
  \item 实现控制平面 Controller，维护 \texttt{ClusterConfig}，支持初始分片分配、MoveShard 和 Rebalance。
  \item 实现配置持久化链路，将控制面配置保存到 etcd，并可附带保存到本地文件，使 API Server 重启后能够恢复最新配置。
  \item 实现分片级迁移锁，在 MoveShard/Rebalance 期间阻止冲突读写，并通过 \texttt{503 + Retry-After} 反馈给客户端。
  \item 实现 API Server 网关，提供 \texttt{/sql}、\texttt{/config}、\texttt{/shards}、\texttt{/groups}、\texttt{/move-shard}、\texttt{/rebalance} 等 HTTP 接口。
  \item 实现 Dashboard 与 CLI 接入，支持拓扑展示、健康探测、分片查看、表数据浏览和控制命令调用。
\end{enumerate}

\docfigure{ddbarch.jpeg}{控制平面与 API 接入在系统架构中的位置}

\section{控制平面元数据管理实现说明}

\subsection{模块组件设计}

控制平面的核心实现位于 \texttt{internal/controller/service.go}。它不直接处理 SQL 语句，而是管理分片配置、分片锁和配置持久化。上层 API Server 通过 Controller Service 获取当前配置、预览迁移动作、更新配置，并在迁移期间持有锁。

本模块包括以下几个部分：
\begin{enumerate}
  \item Controller Service。\texttt{Service} 持有当前 \texttt{ClusterConfig}、配置存储接口和 \texttt{lockedShardIDs} 集合，提供 \texttt{CurrentConfig}、\texttt{PreviewMoveShard}、\texttt{PreviewRebalance}、\texttt{UpdateConfig}、\texttt{LockShards}、\texttt{UnlockShards} 与 \texttt{WithLockedShards} 等方法。
  \item 配置初始化。\texttt{NewBootstrapService} 优先从存储中加载已有配置；若配置不存在，则根据总分片数和启动时发现的 groupIDs 进行轮询式初始分配，并将配置保存。默认总分片数为 8。
  \item 配置存储抽象。\texttt{ConfigStore} 定义 \texttt{Load} 和 \texttt{Save} 两个接口，具体实现包括内存存储、文件存储、基于 etcd 的 DiscoveryStore 以及 ChainStore。ChainStore 可按顺序读取多个存储，并在保存时写入多个后端。
  \item etcd 服务发现。\texttt{internal/discovery/etcd.go} 使用 \texttt{/ddb/nodes/} 保存带 Lease 的在线节点信息，使用 \texttt{/ddb/removed/} 保存被移除节点标记，使用 \texttt{/ddb/controller/config} 保存控制平面配置。节点注册使用 10 秒 Lease 与 KeepAlive，节点离线后注册键自动过期。
  \item 应用装配。\texttt{internal/app/app.go} 根据启动角色选择不同服务：\texttt{controller} 与 \texttt{apiserver} 角色会装配 Controller Service、Router 和 Coordinator；\texttt{shard} 角色会装配 BoltDB、RaftNode 和普通数据节点 API。
  \item 分片迁移编排。MoveShard/Rebalance 并不只修改元数据，而是在 \texttt{WithLockedShards} 保护下先调用 migrator 搬迁数据，再更新 \texttt{ClusterConfig}，最后释放锁。
\end{enumerate}

\subsection{主要数据结构}

控制平面主要数据结构定义在 \texttt{internal/shardmeta/types.go}、\texttt{internal/controller/service.go} 和 \texttt{internal/model/model.go} 中：
\begin{enumerate}
  \item \texttt{ClusterConfig}：全局分片配置，包含版本号、总分片数、分片分配列表和可选副本组信息。版本号在 MoveShard/Rebalance 后递增，用于表示配置演进。
  \item \texttt{ShardAssignment}：单个分片到副本组的映射，包含 \texttt{ShardID} 与 \texttt{GroupID}。系统通过该结构决定某个分片当前由哪个 Raft 副本组承载。
  \item \texttt{GroupInfo}：副本组信息，包含组 ID 和可选节点列表，用于描述数据平面的逻辑承载单元。
  \item \texttt{NodeRole}：节点角色枚举，包括 \texttt{shard}、\texttt{controller}、\texttt{apiserver}。角色化启动使同一个 \texttt{cmd/server} 可以运行不同职责的进程。
  \item \texttt{NodeInfo}：服务发现中的节点信息，包含节点 ID、HTTP 地址、Raft 地址、Leader 状态、角色和 GroupID。API Server 通过该结构挑选目标副本组的 Leader 或在线节点。
  \item \texttt{ShardMigrationError}：迁移锁冲突错误，封装具体分片 ID，并可 unwrap 到 \texttt{ErrShardMigrationInProgress}，便于 API 层转换为 503 响应。
\end{enumerate}

\subsubsection{控制平面类图}

\begin{lstlisting}[style=asciiart,caption={控制平面元数据管理类图}]
+-----------------------------+          +-----------------------------+
| controller.Service          |          | controller.ConfigStore      |
+-----------------------------+          +-----------------------------+
| config: ClusterConfig       | -------> | Load(ctx) ClusterConfig     |
| lockedShardIDs: map         |          | Save(ctx, config) error     |
+-----------------------------+          +-----------------------------+
| CurrentConfig()             |              ^        ^          ^
| PreviewMoveShard(id, group) |              |        |          |
| PreviewRebalance(groups)    |   +----------+   +----+-----+    +----------------+
| UpdateConfig(config)        |   |              |          |                     |
| WithLockedShards(ids, fn)   |   |              |          |                     |
+-----------------------------+   |              |          |                     |
                                  |              |          |                     |
+-----------------------------+   |  +------------------+  +-------------------+
| shardmeta.ClusterConfig     |   |  | MemoryStore      |  | FileStore         |
+-----------------------------+   |  +------------------+  +-------------------+
| Version: uint64             |   |
| TotalShards: int            |   |  +------------------+      +----------------+
| Assignments: []ShardAssign  |   +--| DiscoveryStore   | ---> | discovery.Client|
| Groups: []GroupInfo         |      +------------------+      +----------------+
+-----------------------------+                                | ListNodes()    |
         |                                                     | Register()     |
         | contains                                            | SaveConfig()   |
         v                                                     | LoadConfig()   |
+-----------------------------+                                +----------------+
| shardmeta.ShardAssignment   |
+-----------------------------+
| ShardID: ShardID            |
| GroupID: GroupID            |
+-----------------------------+
\end{lstlisting}

\subsection{流程图设计}

\subsubsection{控制平面启动流程}

控制平面启动时，\texttt{cmd/server} 读取命令行参数和环境变量，生成 \texttt{ServerConfig}。若角色为 \texttt{controller} 或 \texttt{apiserver}，\texttt{NewServerApp} 会连接 etcd，读取当前在线 shard 节点并选择启动分组；随后创建 \texttt{DiscoveryStore + FileStore} 的 ChainStore，并调用 \texttt{NewBootstrapService} 加载或初始化 \texttt{ClusterConfig}。最后创建 Router 与 Coordinator，并将它们挂载到 API Server Handler。

\subsubsection{MoveShard / Rebalance 流程}

MoveShard 流程首先读取当前配置，根据目标 shard 找到源副本组；然后 Controller 对该分片加锁，Coordinator 从源组读取表结构和行数据，将属于该分片的数据逐行写入目标组并从源组删除；迁移完成后 Controller 生成新配置、递增版本号、持久化配置并释放锁。Rebalance 与之类似，只是先根据新 group 列表计算多条分片移动计划，再按计划逐个迁移。

\docfigure{7a26e3c261844642abe172d3d0195046.png}{分片迁移与控制平面锁保护示意图}

\section{API 接入与可观测面实现说明}

\subsection{模块组件设计}

API 接入层分为两类：一类是 shard 数据节点上的普通 DDB API，直接服务于某个 Raft 副本组；另一类是 controller/apiserver 角色上的 ShardDB API Server，用于统一接收 SQL、控制命令和 Dashboard 请求。

本模块包括以下几个部分：
\begin{enumerate}
  \item 数据节点 API。\texttt{internal/api/handler.go} 提供 \texttt{/health}、\texttt{/sql}、\texttt{/status}、\texttt{/leader}、\texttt{/members}、\texttt{/tables}、\texttt{/schema}、\texttt{/join}、\texttt{/remove} 与 \texttt{/rejoin}。这些接口主要围绕单个 Raft 副本组工作。
  \item ShardDB API Server。\texttt{internal/apiserver/handler.go} 提供 \texttt{/sql}、\texttt{/config}、\texttt{/shards}、\texttt{/groups}、\texttt{/move-shard}、\texttt{/rebalance}。它通过 Controller Service 读取元数据，通过 Coordinator 执行 SQL 或迁移。
  \item SQL 接入路径。API Server 收到 \texttt{/sql} 后，将请求交给 Coordinator。Coordinator 解析 SQL 后，根据语句类型选择单分片路由、广播、Scatter-Gather 或 JOIN 聚合策略，再通过 HTTP 调用目标 shard 节点。
  \item 控制命令接入路径。\texttt{/move-shard} 和 \texttt{/rebalance} 先进行参数校验，再调用 Controller 预览配置变更，并在锁保护下执行迁移和配置更新。
  \item Dashboard。\texttt{internal/apiserver/dashboard.go} 使用 \texttt{go:embed} 内嵌静态资源，提供 \texttt{/dashboard/} 页面、\texttt{/dashboard/api/overview} 拓扑概览和 \texttt{/dashboard/api/table-data} 表数据浏览接口。
  \item CLI。\texttt{cmd/cli/main.go} 提供 \texttt{sql}、\texttt{cluster}、\texttt{control} 和 \texttt{interact} 等命令。CLI 可通过 \texttt{--node-url} 指定接入地址，也可通过 etcd 自动发现 API Server、Controller 或 Leader 节点。
\end{enumerate}

\subsection{主要数据结构}

\begin{enumerate}
  \item \texttt{MoveShardRequest}：\texttt{/move-shard} 请求体，包含目标 \texttt{shard\_id} 与新的 \texttt{group\_id}。
  \item \texttt{RebalanceRequest}：\texttt{/rebalance} 请求体，包含新的副本组 ID 列表，Controller 根据该列表重新均衡分片。
  \item \texttt{model.SQLRequest} 与 \texttt{model.SQLResponse}：SQL API 的统一输入输出结构。响应中包含 \texttt{success}、\texttt{result}、\texttt{error} 和可选 \texttt{leader} 字段。
  \item \texttt{model.ShardsResponse} 与 \texttt{model.GroupStatus}：控制面查询接口的输出结构，分别用于展示分片分配和副本组状态。
  \item \texttt{dashboardOverview}：Dashboard 聚合结构，包含生成时间、摘要、当前配置、分片状态、锁定分片、节点状态、组状态和错误列表。
  \item \texttt{dashboardNode} 与 \texttt{dashboardGroup}：Dashboard 展示层结构，额外记录节点可达性、表数量、Leader 信息、分组健康状态和迁移状态。
\end{enumerate}

\subsubsection{API 接入类图}

\begin{lstlisting}[style=asciiart,caption={API 接入与 Dashboard 聚合类图}]
+-----------------------------+        +-----------------------------+
| apiserver.Handler           |        | coordinator.Coordinator     |
+-----------------------------+        +-----------------------------+
| /sql                        | -----> | configReader: ConfigReader  |
| /config                     |        | nodeLister: NodeLister      |
| /shards                     |        | router: *router.Router      |
| /groups                     |        | httpClient: *http.Client    |
| /move-shard                 |        +-----------------------------+
| /rebalance                  |        | ExecuteSQL(ctx, sql)        |
+-----------------------------+        | MigrateShard(ctx, id, g1,g2)|
          |                            +-----------------------------+
          | uses                                      |
          v                                           | reads and routes
+-----------------------------+        +-----------------------------+
| controller.Service          |        | router.Router               |
+-----------------------------+        +-----------------------------+
| CurrentConfig()             |        | Route(table, key, config)   |
| WithLockedShards(ids, fn)   |        +-----------------------------+
+-----------------------------+                     |
          |                                         v
          | returns                         +-----------------------------+
          v                                 | model.SQLResponse           |
+-----------------------------+             +-----------------------------+
| model.ShardsResponse        |             | Success, Result, Error      |
| model.GroupStatus           |             | Leader                      |
+-----------------------------+             +-----------------------------+

+-----------------------------+        +-----------------------------+
| dashboardOverview           |        | dashboardNode/dashboardGroup|
+-----------------------------+        +-----------------------------+
| Summary, Config, Shards     | -----> | health, leader, tables      |
| LockedShards, Nodes, Groups |        | group status and shard list |
+-----------------------------+        +-----------------------------+
\end{lstlisting}

\subsection{流程图设计}

\subsubsection{SQL API 接入流程}

SQL API 的核心流程为：客户端发送 \texttt{POST /sql} 到 API Server；API Server 解码 \texttt{SQLRequest} 并调用 Coordinator；Coordinator 解析 SQL 并读取当前 \texttt{ClusterConfig}；若语句带主键条件，则通过 Router 定位目标 shard 与 group；若是 CREATE TABLE 则广播到所有 group；若是无主键 SELECT 或 JOIN，则 Scatter 到多个 group 后聚合结果；最后 API Server 将统一的 \texttt{SQLResponse} 返回客户端。

\docfigure{78b5598871104db7ab8998a0b3b3bdcf.png}{SQL API 接入与路由流程图}

\subsubsection{Dashboard 数据聚合流程}

Dashboard 概览接口不直接读取数据库全量数据，而是轻量聚合控制面和节点健康信息。\texttt{/dashboard/api/overview} 首先读取当前 \texttt{ClusterConfig} 和锁定分片，再通过 etcd 获取在线节点；随后对每个节点执行 \texttt{/health} 探测，对 shard 节点额外读取 \texttt{/status} 获取 Leader 和表列表；最后按 GroupID 聚合节点与分片，计算 healthy、degraded、offline、migrating 等状态。

表数据浏览接口 \texttt{/dashboard/api/table-data?table=\ldots} 使用 Coordinator 执行 \texttt{SELECT * FROM 表名}，由统一 SQL 路径完成跨分片查询，避免 Dashboard 自己重复实现路由和聚合逻辑。

\section{测试结果}

\subsection{控制平面功能测试}

\subsubsection{测试用例}

控制平面相关测试主要覆盖以下场景：
\begin{enumerate}
  \item \texttt{TestMoveShard}：验证单个分片从源组移动到目标组后，\texttt{ClusterConfig} 中的分片归属正确更新。
  \item \texttt{TestRebalance}：验证给定新的 group 列表后，系统能够重新生成均衡的分片分配。
  \item \texttt{TestSharedStorePropagatesConfigAcrossServices}：验证多个 Controller Service 通过共享存储看到一致的配置更新。
  \item \texttt{TestShardLocks}：验证分片迁移锁能够阻止重复加锁，并在释放后允许后续操作。
  \item \texttt{TestFileStorePersistsConfig} 与 \texttt{TestEmbeddedClientSavesAndLoadsControllerConfig}：验证配置可以持久化到文件和 etcd，并在重启后加载。
\end{enumerate}

\subsubsection{测试结果}

项目测试已全部通过。结果说明控制平面对分片配置的合法性校验、版本递增、共享配置传播、分片锁保护和持久化恢复均符合预期；端到端测试中的 ShardDB move-shard 与 API Server restart 场景也验证了控制面配置在真实进程环境中的可恢复性。

\subsection{API 接入与 Dashboard 测试}

\subsubsection{测试用例}

API 接入相关测试主要覆盖以下场景：
\begin{enumerate}
  \item \texttt{TestNewHandlerHealthConfigAndControlOperations}：验证 \texttt{/health}、\texttt{/config}、\texttt{/shards}、\texttt{/groups}、\texttt{/move-shard} 与 \texttt{/rebalance} 的正常返回。
  \item \texttt{TestNewHandlerControlErrors}：验证非法方法、未知 shard、空 rebalance group 列表等错误路径。
  \item \texttt{TestNewHandlerSQL}：验证 \texttt{/sql} 能够调用统一 SQL Executor。
  \item \texttt{TestNewHandlerSQLReturnsRetryDuringShardMigration}：验证迁移冲突时返回 HTTP 503 并带 \texttt{Retry-After}。
  \item \texttt{TestNewHandlerDashboardRoutes} 与 \texttt{TestNewHandlerDashboardTableData}：验证 Dashboard 静态页面、概览接口和表数据接口。
  \item 端到端测试中的 \texttt{TestShardMoveMigratesRowsAndUpdatesControlPlane}、\texttt{TestScatterSelectAndEqualityJoinAcrossShards}、\texttt{TestShardMoveReturnsRetryDuringMigration}：验证 API、Coordinator、Controller 与数据节点协作路径。
\end{enumerate}

\subsubsection{测试结果}

项目测试已全部通过。测试结果表明，API Server 能够正确处理 SQL 查询、控制面查询、分片迁移、再平衡和 Dashboard 访问；在迁移期间，冲突请求会得到明确的 503 与 Retry-After 反馈，避免客户端误以为写入已经提交。

\section{开发体会}

\begin{enumerate}
  \item 控制平面给我的主要启发是，分布式系统不只需要“能写入数据”，还需要有地方回答“数据现在应该在哪里”。\texttt{ClusterConfig} 看起来只是一个小型元数据结构，但它实际上决定了所有 SQL 请求的路径和迁移后的系统拓扑。
  \item API 接入层的设计更接近用户视角。底层可以有多个 Raft 组、多个 Leader、多个迁移状态，但用户仍然希望通过一个稳定入口发 SQL 或控制命令。因此，API Server 的工作重点不是执行所有逻辑，而是把复杂性收拢到统一协议和统一错误语义里。
  \item 迁移锁的引入让我意识到，控制操作本身也会和业务请求竞争。MoveShard 如果只更新配置而不约束并发写入，很容易出现数据已经搬走但请求还写入旧组的情况；\texttt{503 + Retry-After} 虽然简单，却给客户端提供了清晰的重试边界。
  \item Dashboard 的价值不在于替代命令行，而在于降低观察系统状态的成本。拓扑、分片、锁、节点健康和表数据放在同一个页面后，演示时可以更直观地说明一次路由或迁移背后发生了什么。
  \item 后续如果继续做控制面，我会优先补充两类能力：一是控制操作审计和权限校验，避免误操作；二是更细的迁移进度与失败恢复机制，让 MoveShard 不只是一个同步接口，而是可观察、可恢复的后台任务。
\end{enumerate}

\section*{参考文献}
\addcontentsline{toc}{section}{参考文献}

\begin{enumerate}[label={[\arabic*]}, leftmargin=2.8em, itemsep=0.3em]
  \item etcd contributors. etcd Documentation: A distributed, reliable key-value store. \url{https://etcd.io/docs/}, 2024.
  \item Karger D, Lehman E, Leighton T, et al. Consistent hashing and random trees: Distributed caching protocols for relieving hot spots on the World Wide Web[C]. STOC, 1997.
  \item HashiCorp. HashiCorp Raft: A Go implementation of the Raft consensus protocol. \url{https://github.com/hashicorp/raft}, 2024.
  \item PingCAP. TiKV Documentation: A Distributed Transactional Key-Value Database. \url{https://tikv.org/docs/}, 2024.
\end{enumerate}

\end{document}
